
线性代数

- 矩阵和向量运算：在数据表示和转换中使用，如矩阵乘法、转置和逆矩阵。
- 特征值和特征向量：用于数据降维和特征选择，如主成分分析（PCA）。
- 奇异值分解（SVD）：用于数据压缩、降噪和推荐系统。
- 矩阵分解技术：如非负矩阵分解（NMF）用于主题建模。



微积分

- 导数和偏导数：用于计算函数的变化率，广泛应用于梯度下降算法和神经网络的训练。
- 链式法则：用于反向传播算法，以计算神经网络中的梯度。
- 梯度和Hessian矩阵：用于优化问题，特别是在高维空间中。
- 积分和多重积分：用于概率计算和期望值计算，如在贝叶斯推断中。

概率与统计

- 随机变量和概率分布：用于建模和分析随机现象，如正态分布、泊松分布。
- 贝叶斯定理：用于更新概率分布，广泛应用于贝叶斯网络和贝叶斯推断。
- 最大似然估计（MLE）和最大后验估计（MAP）：用于参数估计和模型选择。
- 马尔可夫链和马尔可夫决策过程（MDP）：用于建模和解决决策过程，特别是在强化学习中。
- 蒙特卡罗方法：用于数值积分和概率分布的近似计算。

优化理论

- 凸优化：研究凸函数和凸集合的优化问题，如LASSO和支持向量机（SVM）。
- 非凸优化：处理复杂的优化问题，如深度学习中的神经网络训练。
- 梯度下降和随机梯度下降：用于迭代优化问题。
- 线性规划和整数规划：用于资源分配和组合优化问题。
- 拉格朗日乘数法：用于含约束条件的优化问题。

图论

- 图的表示：用于建模关系和网络，如邻接矩阵和邻接表。
- 最短路径算法：如Dijkstra算法，用于路径规划和网络分析。
- 图匹配和图划分：用于社交网络分析、图像分割和聚类。
- 网络流和最大流问题：用于优化运输和流量问题。

信息论

- 熵和互信息：用于衡量信息量和依赖关系，在特征选择和通信中应用。
- 编码理论：研究数据压缩和错误检测/纠正码，如香农编码和汉明码。
- 卡尔曼滤波：用于动态系统的状态估计和信号处理。

计算理论

- 图灵机和计算复杂性：研究计算问题的复杂性和可计算性，如P与NP问题。
- NP完全性：研究问题的求解难度和算法的效率，判定问题是否具有多项式时间复杂度解。
- 自动机理论：研究有限状态机和形式语言，应用于编译器和正则表达式。

离散数学

- 组合数学：用于分析和解决离散结构的问题，如组合优化、排列和组合。
- 布尔代数：用于逻辑推理和数理逻辑，特别是在逻辑编程和电路设计中。
- 图论和网络流：用于优化和分析网络结构。

动态系统与微分方程

- 常微分方程（ODEs）：用于建模连续时间的动态系统，如物理模拟和控制理论。
- 偏微分方程（PDEs）：用于建模空间和时间的连续变化，如图像处理和物理场模拟。
- 分岔理论和混沌理论：研究动态系统的稳定性和非线性行为。

几何与拓扑

- 微分几何：用于分析曲面和形状，如在计算机视觉和图像处理中的应用。
- 拓扑数据分析（TDA）：用于数据的形状分析和降维，应用于高维数据的可视化。
- 计算几何：研究几何对象的算法和数据结构，用于计算机图形学和机器人学。

变分法与计算几何

- 变分法：用于优化问题的解，如在图像处理、路径规划和物理模拟中的应用。
- 计算几何：研究几何对象的算法和数据结构，用于图形学、计算机视觉和CAD系统。

数值分析

- 数值微分和积分：用于逼近和求解微分方程和积分方程。
- 线性代数的数值方法：如QR分解、LU分解，用于求解线性方程组。
- 优化的数值方法：如牛顿法、共轭梯度法，用于高效求解优化问题。

笛卡尔几何与代数几何

- 笛卡尔几何：用于解析几何问题，如直线和曲线的表示与计算。
- 代数几何：研究多项式方程的解集结构，用于计算机代数系统和图形学。

逻辑与集合论

- 一阶逻辑和高阶逻辑：用于形式化推理和验证。
- 集合论：研究集合的性质和操作，是数学的基础理论之一。

#2.预备知识

要学习深度学习，首先需要先掌握一些基本技能。所有机器学习方法都涉及从数据中提取信息。因此，我们先学习一些关于数据的实用技能，包括存储、操作和预处理数据。

机器学习通常需要处理大型数据集。我们可以将某些数据集视为一个表，其中表的行对应样本，列对应属性。线性代数为人们提供了一些用来处理表格数据的方法。我们不会太深究细节，而是将重点放在矩阵运算的基本原理及其实现上。

深度学习是关于优化的学习。对于一个带有参数的模型，我们想要找到其中能拟合数据的最好模型。在算法的每个步骤中，决定以何种方式调整参数需要一点微积分知识。本章将简要介绍这些知识。幸运的是，`autograd`包会自动计算微分，本章也将介绍它。

机器学习还涉及如何做出预测：给定观察到的信息，某些未知属性可能的值是多少？要在不确定的情况下进行严格的推断，我们需要借用概率语言。

最后，官方文档提供了本书之外的大量描述和示例。在本章的结尾，我们将展示如何在官方文档中查找所需信息。

本书对读者数学基础无过分要求，只要可以正确理解深度学习所需的数学知识即可。但这并不意味着本书中不涉及数学方面的内容，本章会快速介绍一些基本且常用的数学知识，以便读者能够理解书中的大部分数学内容。如果读者想要深入理解全部数学内容，可以进一步学习本书数学附录中给出的数学基础知识。

-2.1.数据操作
-[2.1.1.入门](https：//zh-v2.d2l.ai/chapter_preliminaries/ndarray.html#id2)
-[2.1.2.运算符](https：//zh-v2.d2l.ai/chapter_preliminaries/ndarray.html#id3)
-[2.1.3.广播机制](https：//zh-v2.d2l.ai/chapter_preliminaries/ndarray.html#subsec-broadcasting)
-[2.1.4.索引和切片](https：//zh-v2.d2l.ai/chapter_preliminaries/ndarray.html#id5)
-[2.1.5.节省内存](https：//zh-v2.d2l.ai/chapter_preliminaries/ndarray.html#id6)
-[2.1.6.转换为其他Python对象](https：//zh-v2.d2l.ai/chapter_preliminaries/ndarray.html#python)
-[2.1.7.小结](https：//zh-v2.d2l.ai/chapter_preliminaries/ndarray.html#id7)
-[2.1.8.练习](https：//zh-v2.d2l.ai/chapter_preliminaries/ndarray.html#id8)
-2.2.数据预处理
-[2.2.1.读取数据集](https：//zh-v2.d2l.ai/chapter_preliminaries/pandas.html#id2)
-[2.2.2.处理缺失值](https：//zh-v2.d2l.ai/chapter_preliminaries/pandas.html#id3)
-[2.2.3.转换为张量格式](https：//zh-v2.d2l.ai/chapter_preliminaries/pandas.html#id4)
-[2.2.4.小结](https：//zh-v2.d2l.ai/chapter_preliminaries/pandas.html#id5)
-[2.2.5.练习](https：//zh-v2.d2l.ai/chapter_preliminaries/pandas.html#id6)
-2.3.线性代数
-[2.3.1.标量](https：//zh-v2.d2l.ai/chapter_preliminaries/linear-algebra.html#id2)
-[2.3.2.向量](https：//zh-v2.d2l.ai/chapter_preliminaries/linear-algebra.html#id3)
-[2.3.3.矩阵](https：//zh-v2.d2l.ai/chapter_preliminaries/linear-algebra.html#id5)
-[2.3.4.张量](https：//zh-v2.d2l.ai/chapter_preliminaries/linear-algebra.html#id6)
-[2.3.5.张量算法的基本性质](https：//zh-v2.d2l.ai/chapter_preliminaries/linear-algebra.html#id7)
-[2.3.6.降维](https：//zh-v2.d2l.ai/chapter_preliminaries/linear-algebra.html#subseq-lin-alg-reduction)
-[2.3.7.点积（DotProduct）](https：//zh-v2.d2l.ai/chapter_preliminaries/linear-algebra.html#dot-product)
-[2.3.8.矩阵-向量积](https：//zh-v2.d2l.ai/chapter_preliminaries/linear-algebra.html#id10)
-[2.3.9.矩阵-矩阵乘法](https：//zh-v2.d2l.ai/chapter_preliminaries/linear-algebra.html#id11)
-[2.3.10.范数](https：//zh-v2.d2l.ai/chapter_preliminaries/linear-algebra.html#subsec-lin-algebra-norms)
-[2.3.11.关于线性代数的更多信息](https：//zh-v2.d2l.ai/chapter_preliminaries/linear-algebra.html#id14)
-[2.3.12.小结](https：//zh-v2.d2l.ai/chapter_preliminaries/linear-algebra.html#id16)
-[2.3.13.练习](https：//zh-v2.d2l.ai/chapter_preliminaries/linear-algebra.html#id17)
-2.4.微积分
-[2.4.1.导数和微分](https：//zh-v2.d2l.ai/chapter_preliminaries/calculus.html#id2)
-[2.4.2.偏导数](https：//zh-v2.d2l.ai/chapter_preliminaries/calculus.html#id3)
-[2.4.3.梯度](https：//zh-v2.d2l.ai/chapter_preliminaries/calculus.html#subsec-calculus-grad)
-[2.4.4.链式法则](https：//zh-v2.d2l.ai/chapter_preliminaries/calculus.html#id5)
-[2.4.5.小结](https：//zh-v2.d2l.ai/chapter_preliminaries/calculus.html#id6)
-[2.4.6.练习](https：//zh-v2.d2l.ai/chapter_preliminaries/calculus.html#id7)
-2.5.自动微分
-[2.5.1.一个简单的例子](https：//zh-v2.d2l.ai/chapter_preliminaries/autograd.html#id2)
-[2.5.2.非标量变量的反向传播](https：//zh-v2.d2l.ai/chapter_preliminaries/autograd.html#id3)
-[2.5.3.分离计算](https：//zh-v2.d2l.ai/chapter_preliminaries/autograd.html#id4)
-[2.5.4.Python控制流的梯度计算](https：//zh-v2.d2l.ai/chapter_preliminaries/autograd.html#python)
-[2.5.5.小结](https：//zh-v2.d2l.ai/chapter_preliminaries/autograd.html#id5)
-[2.5.6.练习](https：//zh-v2.d2l.ai/chapter_preliminaries/autograd.html#id6)
-2.6.概率
-[2.6.1.基本概率论](https：//zh-v2.d2l.ai/chapter_preliminaries/probability.html#id2)
-[2.6.2.处理多个随机变量](https：//zh-v2.d2l.ai/chapter_preliminaries/probability.html#id5)
-[2.6.3.期望和方差](https：//zh-v2.d2l.ai/chapter_preliminaries/probability.html#id12)
-[2.6.4.小结](https：//zh-v2.d2l.ai/chapter_preliminaries/probability.html#id13)
-[2.6.5.练习](https：//zh-v2.d2l.ai/chapter_preliminaries/probability.html#id14)
-2.7.查阅文档
-[2.7.1.查找模块中的所有函数和类](https：//zh-v2.d2l.ai/chapter_preliminaries/lookup-api.html#id2)
-[2.7.2.查找特定函数和类的用法](https：//zh-v2.d2l.ai/chapter_preliminaries/lookup-api.html#id3)
-[2.7.3.小结](https：//zh-v2.d2l.ai/chapter_preliminaries/lookup-api.html#id8)
-[2.7.4.练习](https：//zh-v2.d2l.ai/chapter_preliminaries/lookup-api.html#id9)

